% GENERATED FILE. Edit canonical lessons, then run npm run generate:books.
% canonical-source-hash: fnv1a64:ae7f765e50208966
% canonical-lessons: ZH-C496-beishu, ZH-C496-gaoji, ZH-W496-ji-grade, ZH-C496-zhongji, ZH-C496-chuji, ZH-C496-liuji

\chapter{Levels}
\label{ch:zh-levels}
\hlchaptermodality{496}

\begin{chapteropening}
\textbf{By the end of this chapter:} \emph{I can talk about learning by heart, advanced, intermediate and elementary levels, and repeating a year.}

Complete the last lesson of chapter 496: I can talk about learning by heart, advanced, intermediate and elementary levels, and repeating a year.
\end{chapteropening}

\section[bèishū]{bèishū — to learn a text by heart; to recite from memory}

\label{lesson:ZH-C496-beishu}



\begin{quote}
Before the new one: say the Mandarin for to cheat in an exam, then the Mandarin for a hard problem.
\end{quote}

\subsection*{You'll want to know: bèishū}
\emph{bèishū} — \textquotedblleft{}to learn a text by heart; to recite from memory\textquotedblright{}.

\begin{quote}
\emph{kǎoshì yǐqián wǒ měi tiān zǎoshang bèishū} — Before exams I learn my texts by heart every morning.
\end{quote}

Its characters come later; for now it is a word for the ear.

\begin{cousinweb}
\emph{bèishū} is built from \emph{bèi} 'the back' and \zh{书} 'a book', both inherited from Middle Chinese, \emph{pwojH} and \emph{syo} (Guangyun, in Baxter's transcription), which glosses the first 'the back of the body' and also lists it as \emph{bwojH}, 'to turn one's back on'. How the pair came to mean 'to recite from memory' is not settled.
\end{cousinweb}

\begin{grammarlens}[title={What you've built}]
One more word for saying how it was.
\end{grammarlens}

\subsection*{Your turn}
\begin{itemize}
\raggedright
  \item \emph{Say it:} \emph{bèishū}
  \item \emph{Say it:} \emph{bèishū}, once more
  \item \emph{Say it:} \emph{kǎoshì yǐqián wǒ měi tiān zǎoshang bèishū} — Before exams I learn my texts by heart every morning.
  \item \emph{Recall:} say the Mandarin for to cheat in an exam, then the Mandarin for a hard problem, then say \emph{bèishū} again
\end{itemize}

\begin{tcolorbox}[breakable,colback=teal!4,colframe=teal!35!black,title={Before you move on}]
What does bèishū mean? (\textquotedblleft{}To learn a text by heart; to recite from memory\textquotedblright{}.) Say it once more.
\end{tcolorbox}
\section[\texorpdfstring{\zh{高级} (gāojí)}{高级 (gāojí)}]{\zh{高级} (gāojí) — advanced; high-level}

\label{lesson:ZH-C496-gaoji}



\begin{quote}
Before the new one: say the Mandarin for a hard problem, then the Mandarin for to learn a text by heart.
\end{quote}

\subsection*{You'll want to know: \zh{高级}}
\textbf{\zh{高级}} — \emph{gāojí} — \textquotedblleft{}advanced; high-level\textquotedblright{}.

\begin{quote}
\textbf{\zh{我学了三年汉语现在学高级汉语}} \emph{wǒ xué le sān nián Hànyǔ xiànzài xué gāojí Hànyǔ} — I've studied Chinese for three years, and now I'm doing advanced Chinese.
\end{quote}

The new character, in pieces: \zh{级} — six strokes in two pieces: on the left the silk piece of \zh{红} \emph{hóng}, three strokes; on the right three strokes, a long stroke falling left, a stroke that runs right and turns twice, and a long stroke falling right. The next lesson puts it in your hand.

\begin{cousinweb}
\textbf{\zh{高级}} (\emph{gāojí}) is built from \zh{高} 'high' and \zh{级} 'a grade', both inherited from Middle Chinese, \emph{kaw} and \emph{kip} (Guangyun, in Baxter's transcription), which glosses the second 'a grade' and, after an older dictionary, 'the order of silk threads'.
\end{cousinweb}

\begin{grammarlens}[title={What you've built}]
One more word for saying how it was.
\end{grammarlens}

\subsection*{Your turn}
\begin{itemize}
\raggedright
  \item \emph{Say it:} \emph{gāojí}
  \item \emph{Say it:} \emph{gāojí}, once more
  \item \emph{Say it:} \emph{wǒ xué le sān nián Hànyǔ xiànzài xué gāojí Hànyǔ} — I've studied Chinese for three years, and now I'm doing advanced Chinese.
  \item \emph{Recall:} say the Mandarin for a hard problem, then the Mandarin for to learn a text by heart, then say \emph{gāojí} again
\end{itemize}

\begin{tcolorbox}[breakable,colback=teal!4,colframe=teal!35!black,title={Before you move on}]
What does \zh{高级} mean? (\textquotedblleft{}Advanced; high-level\textquotedblright{}.) Say it once more.
\end{tcolorbox}
\section[\texorpdfstring{\zh{级} (jí)}{级 (jí)}]{\zh{级} — the silk piece, then three strokes}

\label{lesson:ZH-W496-ji-grade}



\begin{quote}
The word first, then the character inside it.

\begin{itemize}
\raggedright
  \item \emph{Recall:} say \emph{gāojí} — \textbf{R1}, one lesson back, before it decays
\end{itemize}

\begin{quote}
\zh{红} \zh{高}
\end{quote}

\emph{Read it:} each one aloud, with nothing said first — \emph{hóng}, \emph{gāo}
\end{quote}

\subsection*{Script — the silk piece, then three strokes}
\hlblockfigure{\detokenize{figures/ZH-W496-ji-grade-filmstrip.pdf}}{How \zh{级} is written, stroke by stroke}

\begin{quote}
\zh{级}
\end{quote}

\textbf{Two pieces, left before right.}

\textbf{The left piece, three strokes}, as on the left of \zh{红}: a stroke that falls to the left and turns right, a second one like it below, then a short stroke rising to the right.

\textbf{The right piece, three strokes}: a long stroke falling to the left; then one stroke that runs right along the top, turns down and back, runs right again and falls away to the lower left; then a long stroke falling to the right.

3 + 3 = six.

\subsection*{Your turn}
\begin{itemize}
\raggedright
  \item \emph{Write it:} \zh{级}, all six, with the model in front of you
  \item \emph{Cover:} write it again from memory, and count the pieces, not the strokes
\end{itemize}

\begin{tcolorbox}[breakable,colback=teal!4,colframe=teal!35!black,title={Before you move on}]
How many pieces does \zh{级} have? (\textbf{Two}.) Which characters you know share its left piece? (\textbf{\zh{红}, \zh{经} or \zh{约}}.) What is the last stroke? (\textbf{A long stroke falling to the right}.)

Source: \href{https://raw.githubusercontent.com/chanind/hanzi-writer-data/68d10a4b21150cae5e1ebbd223eed289cf32d90c/data/\%E7\%BA\%A7.json}{Hanzi Writer Data \zh{级}, pinned stroke paths and medians}.
\end{tcolorbox}
\section[\texorpdfstring{\zh{中级} (zhōngjí)}{中级 (zhōngjí)}]{\zh{中级} (zhōngjí) — intermediate (in level)}

\label{lesson:ZH-C496-zhongji}



\begin{quote}
Before the new one: say the Mandarin for to learn a text by heart, then the Mandarin for advanced (in level).
\end{quote}

\subsection*{You'll want to know: \zh{中级}}
\textbf{\zh{中级}} — \emph{zhōngjí} — \textquotedblleft{}intermediate (in level)\textquotedblright{}.

\begin{quote}
\textbf{\zh{我在公司学中级汉语}} \emph{wǒ zài gōngsī xué zhōngjí Hànyǔ} — I'm taking intermediate Chinese at work.
\end{quote}

\begin{cousinweb}
\textbf{\zh{中级}} (\emph{zhōngjí}) is built from \zh{中} 'the middle' and \zh{级} 'a grade', both inherited from Middle Chinese, \emph{trjuwng} and \emph{kip} (Guangyun, in Baxter's transcription).
\end{cousinweb}

\begin{grammarlens}[title={What you've built}]
One more word for saying how it was.
\end{grammarlens}

\subsection*{Your turn}
\begin{itemize}
\raggedright
  \item \emph{Say it:} \emph{zhōngjí}
  \item \emph{Say it:} \emph{zhōngjí}, once more
  \item \emph{Say it:} \emph{wǒ zài gōngsī xué zhōngjí Hànyǔ} — I'm taking intermediate Chinese at work.
  \item \emph{Recall:} say the Mandarin for to learn a text by heart, then the Mandarin for advanced (in level), then say \emph{zhōngjí} again
\end{itemize}

\begin{tcolorbox}[breakable,colback=teal!4,colframe=teal!35!black,title={Before you move on}]
What does \zh{中级} mean? (\textquotedblleft{}Intermediate (in level)\textquotedblright{}.) Say it once more.
\end{tcolorbox}
\section[chūjí]{chūjí — elementary; for beginners}

\label{lesson:ZH-C496-chuji}



\begin{quote}
Before the new one: say the Mandarin for advanced (in level), then the Mandarin for intermediate (in level).
\end{quote}

\subsection*{You'll want to know: chūjí}
\emph{chūjí} — \textquotedblleft{}elementary; for beginners\textquotedblright{}.

\begin{quote}
\emph{wǒ xué guo chūjí Hànyǔ} — I've done elementary Chinese.
\end{quote}

Its characters come later; for now it is a word for the ear.

\begin{cousinweb}
\emph{chūjí} is built from \emph{chū} 'the start' and \zh{级} 'a grade', both inherited from Middle Chinese, \emph{tsrhjo} and \emph{kip} (Guangyun, in Baxter's transcription), which glosses the first 'the beginning'.
\end{cousinweb}

\begin{grammarlens}[title={What you've built}]
One more word for saying how it was.
\end{grammarlens}

\subsection*{Your turn}
\begin{itemize}
\raggedright
  \item \emph{Say it:} \emph{chūjí}
  \item \emph{Say it:} \emph{chūjí}, once more
  \item \emph{Say it:} \emph{wǒ xué guo chūjí Hànyǔ} — I've done elementary Chinese.
  \item \emph{Recall:} say the Mandarin for advanced (in level), then the Mandarin for intermediate (in level), then say \emph{chūjí} again
\end{itemize}

\begin{tcolorbox}[breakable,colback=teal!4,colframe=teal!35!black,title={Before you move on}]
What does chūjí mean? (\textquotedblleft{}Elementary; for beginners\textquotedblright{}.) Say it once more.
\end{tcolorbox}
\section[liújí]{liújí — to repeat a school year}

\label{lesson:ZH-C496-liuji}



\begin{quote}
Before the new one: say the Mandarin for intermediate (in level), then the Mandarin for elementary (in level).
\end{quote}

\subsection*{You'll want to know: liújí}
\emph{liújí} — \textquotedblleft{}to repeat a school year\textquotedblright{}.

\begin{quote}
\emph{tā de chéngjì bù hǎo liújí le} — His marks were poor, so he had to repeat the year.
\end{quote}

Its characters come later; for now it is a word for the ear.

\begin{cousinweb}
\emph{liújí} is built from \emph{liú} 'to stay' and \zh{级} 'a grade', both inherited from Middle Chinese, \emph{ljuw} and \emph{kip} (Guangyun, in Baxter's transcription), which glosses the first 'to stay' and 'to stop'.
\end{cousinweb}

\begin{grammarlens}[title={What you've built}]
One more word for saying how it was.
\end{grammarlens}

\subsection*{Your turn}
\begin{itemize}
\raggedright
  \item \emph{Say it:} \emph{liújí}
  \item \emph{Say it:} \emph{liújí}, once more
  \item \emph{Say it:} \emph{tā de chéngjì bù hǎo liújí le} — His marks were poor, so he had to repeat the year.
  \item \emph{Recall:} say the Mandarin for intermediate (in level), then the Mandarin for elementary (in level), then say \emph{liújí} again
\end{itemize}

\begin{tcolorbox}[breakable,colback=teal!4,colframe=teal!35!black,title={Before you move on}]
What does liújí mean? (\textquotedblleft{}To repeat a school year\textquotedblright{}.) Say it once more.
\end{tcolorbox}
